\chapter{Происхождение физического опорного масштаба \texorpdfstring{$\mu$}{mu} из родительского функционала}
\label{ch:v9-reference-scale-mu-parent-origin}

\section{Контракт абсолютного масштаба}

После Gaussian reference-state запретительный результат физическое переоткрытие требует
величину $\mu$ с пятью свойствами:
\begin{equation}
 \begin{aligned}
 C_1&=\text{energy dimension},&
 C_2&=\text{internal selector},\\
 C_3&=\text{typed map},&
 C_4&=\text{orbit breaking},\\
 C_5&=\text{target independence}.&&
 \end{aligned}
 \label{eq:v9-mu-origin-contract}
\end{equation}
Кандидат проходит только при произведении всех пяти indicators, равном
единице.

\section{Восемь кандидатов корпуса}

Проверены inverse KMS temperature, clock energy, reservoir cutoff,
inverse compactification radius, spectral Dirac scale, observed mass,
vacuum Hessian gap и dimensional transmutation. Их exact indicator matrix
равна
\begin{equation}
 M_\mu=
 \begin{pmatrix}
 1&0&0&0&1\\
 1&0&0&0&1\\
 1&0&0&0&1\\
 1&0&0&0&1\\
 1&0&0&0&1\\
 1&0&0&1&0\\
 0&1&0&0&1\\
 1&0&0&1&1
 \end{pmatrix},
 \qquad \rank M_\mu=4.
 \label{eq:v9-mu-origin-candidate-matrix}
\end{equation}
Третий столбец нулевой: ни один кандидат не имеет выведенного typed map в
Gaussian covariance carrier. Полный pass-vector поэтому равен
\begin{equation}
 \mathbf p_\mu=
 \left(\prod_{a=1}^{5}M_{ja}\right)_{j=1}^{8}
 =(0,0,0,0,0,0,0,0)^T.
 \label{eq:v9-mu-origin-pass-vector}
\end{equation}

Первые пять кандидатов имеют правильную размерность, но их absolute value
не выбран. Observed mass разрывает orbit только как содержащий целевую подстановку datum.
Vacuum Hessian internally вычислим, но безразмерен. Dimensional
transmutation достигает score $3/5$ и является ближайшим классом,
\begin{equation}
 \mu_{\rm DT}=\Lambda
 \exp\!\left(-\int^{g(\Lambda)}\frac{dg}{\beta(g)}\right),
 \label{eq:v9-mu-origin-dimensional-transmutation}
\end{equation}
однако current parent не выводит ни ненулевую $\beta$-function этого
carrier, ни boundary datum $g(\Lambda)$, ни typed map в $Q$.

\section{Общий scale-orbit}

Пусть семь dimensionful candidates сравниваются с $\mu$. В logarithmic
coordinates карта всех отношений имеет вид
\begin{equation}
 A_{\rm rel}=
 \begin{pmatrix}
 -1&1&0&0&0&0&0&0\\
 -1&0&1&0&0&0&0&0\\
 -1&0&0&1&0&0&0&0\\
 -1&0&0&0&1&0&0&0\\
 -1&0&0&0&0&1&0&0\\
 -1&0&0&0&0&0&1&0\\
 -1&0&0&0&0&0&0&1
 \end{pmatrix}.
 \label{eq:v9-mu-origin-relative-map}
\end{equation}
Точно
\begin{equation}
 \rank A_{\rm rel}=7,
 \qquad \operatorname{nullity}A_{\rm rel}=1,
 \qquad A_{\rm rel}(1,1,1,1,1,1,1,1)^T=0.
 \label{eq:v9-mu-origin-relative-rank}
\end{equation}
Следовательно, сколько угодно relative calibrations не фиксируют один
common absolute energy scale.

Это видно на известных invariants:
\begin{equation}
 \beta E_*,\qquad \frac{E_C\tau_C}{\hbar},\qquad
 \frac{\mu R}{\hbar c}.
 \label{eq:v9-mu-origin-known-invariants}
\end{equation}
Они сохраняются при согласованном rescaling energy и reciprocal time/length.
Repeated-interaction limits действительно связывают coupling и time-step,
но не выбирают абсолютную единицу без microscopic energy input
\cite{v9-mu-origin-attal2006}.

\begin{theorem}[Исчерпание reference-scale кандидатов]
Ни один из восьми кандидатов текущего корпуса не проходит пятичастный
physical-origin contract для $\mu$. Relative-scale map имеет единственную
common scaling zero mode.
\end{theorem}

\begin{nogocriterion}[Абсолютный масштаб не возникает из отношения]
KMS temperature, clock frequency, cutoff, radius или mass могут
калибровать $\mu$ только после independently выбранного dimensionful datum.
Observed mass запрещена как содержащий целевую подстановку input. Dimensional transmutation
переоткрывает путь лишь после вывода новой $\beta$-function и её boundary
condition из microscopic parent.
\end{nogocriterion}

\begin{protocol}\begin{enumerate}
\item candidates: \textbf{$8$};
\item origin criteria: \textbf{$5$};
\item candidate-matrix rank: \textbf{$4$};
\item passing candidates: \textbf{$0/8$};
\item relative-scale rank/nullity: \textbf{$7/1$};
\item physical опорный масштаб $\mu$: \textbf{$0/1$};
\item physical-origin packages: \textbf{$0/2$};
\item Tome IX final physical запретительный результат: \textbf{готов к фиксации}.
\end{enumerate}\end{protocol}

Машинный сертификат сохранён в
\path{s2t_v9_physical_reopening_reference_scale_mu_parent_origin_gate_results.json}.

\begin{thebibliography}{9}
\bibitem{v9-mu-origin-attal2006}
S.~Attal, Y.~Pautrat,
\emph{From repeated to continuous quantum interactions},
Annales Henri Poincar\'e 7 (2006), 59--104; arXiv:math-ph/0311002.
\end{thebibliography}